\documentclass{article}

% NeurIPS 2026 style (default = main track, double-blind)
\usepackage{neurips_2026}

% Common packages
\usepackage[utf8]{inputenc}
\usepackage[T1]{fontenc}
\usepackage{hyperref}
\usepackage{url}
\usepackage{booktabs}
\usepackage{amsfonts}
\usepackage{nicefrac}
\usepackage{microtype}
\usepackage{xcolor}

% Additional packages from the ICML version
\usepackage{graphicx}
\usepackage{subcaption}
\usepackage{multirow}
\usepackage{tabularx}
\usepackage{threeparttable}
\usepackage{amsmath}
\usepackage{amssymb}
\usepackage{mathtools}
\usepackage{amsthm}
\usepackage{pifont}

% Plotting
\usepackage{tikz}
\usepackage{pgfplots}
\pgfplotsset{compat=1.18}

% cleveref (load after hyperref)
\usepackage[capitalize,noabbrev]{cleveref}

%%%%%%%%%%%%%%%%%%%%%%%%%%%%%%%%
% THEOREMS
%%%%%%%%%%%%%%%%%%%%%%%%%%%%%%%%
\theoremstyle{plain}
\newtheorem{theorem}{Theorem}[section]
\newtheorem{proposition}[theorem]{Proposition}
\newtheorem{lemma}[theorem]{Lemma}
\newtheorem{corollary}[theorem]{Corollary}
\theoremstyle{definition}
\newtheorem{definition}[theorem]{Definition}
\newtheorem{assumption}[theorem]{Assumption}
\theoremstyle{remark}
\newtheorem{remark}[theorem]{Remark}

\title{LASER: Lossless ANN-to-SNN Conversion with Error-bounded Nonlinear Adaptation}

\author{%
  Zhengzheng Tang \\
  Department of Computer Science\\
  Boston University\\
  Boston, MA, USA \\
  \texttt{zztangbu@bu.edu} \\
}

\begin{document}

\maketitle

\begin{abstract}
Deploying large language models on energy-efficient neuromorphic hardware requires
accurate ANN-to-SNN conversion, yet existing methods degrade model quality
catastrophically at scale. We identify the root cause as a misclassification of the
problem: prior work treats conversion as an approximation task, accepting lossy spike
encoding at every layer and relying on surrogate gradients to compensate, causing
errors to compound exponentially with network depth. We propose \textbf{LASER}, which
reframes conversion as a representation problem. Bit Spike Encoding (BSE) establishes
a bijective mapping between $N$-bit values and $N$-timestep spike trains, making all
linear computations provably equivalent to their ANN counterparts and eliminating
approximation error from the linear path entirely. Adaptive Spiking Neural Codec
(ASNC) approximates nonlinear activations via distribution-aware piecewise adaptation
without architectural modification, strictly confining the unavoidable approximation
error to a single bounded module. We prove that under BSE and ASNC, the
Straight-Through Estimator is exact for linear layers by construction and
bounded-error at nonlinear layers, replacing the exponential error accumulation of
prior methods with a bound that is linear in the number of nonlinear modules and
independent of network depth. Empirically, LASER achieves near-lossless conversion of
billion-scale language models with accuracy degradation well below that of prior SNN
methods, while reducing nonlinear computation energy by two orders of magnitude on
neuromorphic hardware.
\end{abstract}

\section{Introduction}
\label{sec:introduction}

The rapid scaling of large language models (LLMs) has delivered remarkable capabilities, but at substantial inference cost. Serving a 70-billion-parameter model requires dense GPU clusters consuming kilowatts of power per query—a deployment constraint that grows more acute as models scale further~\citep{tay2022efficient}. Neuromorphic hardware, exemplified by Intel's Loihi~2, offers a fundamentally different computational substrate: event-driven, sparse, and asynchronous, with energy consumption that scales with neural activity rather than model size~\citep{davies2018loihi,davies2021loihi}. Spiking Neural Networks (SNNs),
which communicate via discrete binary spikes, are the natural programming model for such hardware~\citep{maass1997networks,roy2019towards}. If a pretrained LLM could be accurately converted to an SNN, it would become deployable on neuromorphic hardware—preserving the capability earned through months of pretraining while dramatically reducing inference energy.

Existing ANN-to-SNN conversion methods introduce approximation at every stage of the
network: during spike encoding, during nonlinear computation, and during training via
surrogate gradients~\citep{neftci2019surrogate,wu2018stbp,rueckauer2017conversion}.
Each approximation injects error, and in a deep network these errors compound layer by
layer. For small image classifiers with tens of layers this degradation is manageable.
For large language models with hundreds of transformer layers it is catastrophic---prior SNN methods applied to billion-parameter models yield perplexity degradations so severe that the converted model becomes unusable for any practical task~\citep{xing2025spikellm}.

The failure is structural, not incidental. At the encoding stage, converting continuous activations to discrete spikes is inherently lossy: rate coding requires thousands of timesteps to approximate a single value with acceptable
precision~\citep{rueckauer2017conversion,auge2021encoding}, while time-to-first-spike
coding demands orders-of-magnitude more timesteps to approach machine precision~\citep{bonilla2022ttfs,stanojevic2024ttfs}, as shown in
Table~\ref{tab:bse_fidelity}. At the nonlinear stage, functions such as SiLU and
Softmax have no exact spike-domain equivalent, so prior works replace them with
spike-friendly surrogates such as
ReLU~\citep{zhou2023spikformer,yao2023sdt,spikedattention2024}, altering the network
architecture and discarding the carefully optimized behavior of the pretrained model.
At the training stage, the non-differentiability of spikes forces reliance on surrogate gradients~\citep{neftci2019surrogate,shrestha2018slayer,deng2022tet}, which inject additional approximation errors that diffuse through the network and destabilize training at scale. These three error sources interact and accumulate across depth, making large-scale conversion fundamentally brittle under existing approaches.

Prior work treats ANN-to-SNN conversion as an \emph{approximation} problem---represent
continuous values with discrete spikes as accurately as possible, accept inevitable loss, and compensate with longer timesteps or architectural substitution. The correct framing is a \emph{representation} problem: establish an exact, invertible correspondence between the spike domain and the value domain so that linear computation in the spiking network is provably equivalent to its ANN counterpart. Once this correspondence exists, approximation error is not a diffuse, layer-wise phenomenon---it becomes structurally confined to the one location where it is genuinely unavoidable: nonlinear activation functions.
%
We propose \textbf{LASER}, a high-fidelity ANN-to-SNN conversion framework built on
this principle. \textbf{Bit Spike Encoding (BSE)} establishes a bijective mapping
between numerical values and binary spike trains, making all linear computations in
the spike domain provably equivalent to their ANN counterparts---eliminating
approximation error from the linear path entirely. \textbf{Adaptive Spiking Neural
Codec (ASNC)} handles nonlinear activations as a universal piecewise approximator,
fitting arbitrary functions in the spike domain without modifying the network
architecture, strictly confining the unavoidable approximation error to a single
localized module. Built on this high-fidelity forward path, \textbf{STE training}
is no longer a heuristic surrogate but a principled procedure---exact where the
forward pass is exact, and bounded-error where it is not. Our main contribution is:
\begin{itemize}
  \item We propose \textbf{BSE}, an exact bijective encoding between $N$-bit values and $N$-timestep spike trains, proving that all linear computations under BSE are mathematically equivalent to their ANN counterparts, eliminating approximation error from the linear path entirely.

  \item We propose \textbf{ASNC}, a universal nonlinear approximator that fits arbitrary activation functions in the spike domain via distribution-aware piecewise adaptation, without architectural modification, with a provably tighter error bound under Gaussian inputs than prior methods.

  \item We show that under BSE and ASNC, \textbf{STE training is principled rather than heuristic}: exact for linear layers by construction and bounded-error at nonlinear layers by the ASNC guarantee—providing the first rigorous theoretical justification for STE in the SNN setting.

  \item We demonstrate \textbf{near-lossless conversion of 70B-scale LLMs to SNNs}, with less than $2\%$ accuracy degradation across four benchmarks, using only MMLU calibration data and without modifying pretrained weights.
\end{itemize}

\section{Related Work}
\label{sec:related_works}

ANN-to-SNN conversion has been studied along three axes: spike encoding, nonlinear
computation, and training, each of which contributes its own approximation errors.
A common thread across all three is that existing methods accept lossy approximation
as unavoidable, and attempt to compensate through additional compute or architectural
change. We discuss each axis in turn and identify the shared structural limitation.

\noindent\textbf{Spike encoding.}
The dominant encoding schemes, rate coding and time-to-first-spike
(TTFS), approximate continuous values by accumulating spike statistics over a time
window~\citep{rueckauer2017conversion,auge2021encoding,guo2021coding}. Rate coding
requires thousands of timesteps to represent a single value with acceptable precision,
amplifying jitter and making large-scale deployment impractical. TTFS reduces latency
at the cost of a single-spike constraint that severely limits representational
capacity~\citep{bonilla2022ttfs,stanojevic2024ttfs}; TTFS-Former, despite recent
progress, still requires $1{,}024$ steps to approach FP16
precision~\citep{ZhaoHuangDingYu2025_TTFSFormer}. Methods such as
Spikformer/Spikformer-v2~\citep{zhou2023spikformer,zhou2024spikformerv2} and
SpikeZIP/SpikeZIP-TF~\citep{you2024spikeziptf} compensate by using long time windows
or high firing rates, trading latency for accuracy. Across all these schemes, encoding
is a lossy, many-to-one mapping, and this is a fundamental source of error that no
downstream correction can fully recover.

\noindent\textbf{Nonlinear computation.}
Nonlinear functions such as SiLU, Softmax, and LayerNorm have no exact spike-domain
equivalent, forcing one of two compromises. The first is structural substitution:
replace the original nonlinearity with a spike-friendly surrogate, such as the
Spiking Self-Attention (SSA) used in Spikformer~\citep{zhou2023spikformer} or the
structural rewrites in Spike-Driven
Transformer~\citep{yao2023sdt,spikedattention2024,wang2023stsa}. This discards the
carefully optimized behavior of the pretrained model and requires retraining for
alignment. The second is approximate equivalence: provide spike-domain approximations
of Softmax and LayerNorm~\citep{you2024spikeziptf,tang2024sorbet}, or adopt soft
gating and approximately differentiable units as in SpikeGPT~\citep{zhu2023spikegpt}
and Spiking-LLM~\citep{xing2025spikellm}. Both paths inject approximation errors at
every nonlinear site, with no mechanism to prevent those errors from propagating
through subsequent layers.

\noindent\textbf{Training.}
The non-differentiability of the spike function makes standard backpropagation
inapplicable, and the field has converged on surrogate gradient
methods~\citep{neftci2019surrogate,wu2018stbp,shrestha2018slayer} as the default
remedy. STBP~\citep{wu2018stbp}, SLAYER~\citep{shrestha2018slayer}, and
DIET-SNN~\citep{rathi2021dietsnn} introduce surrogate gradients at multiple levels of
the network, reducing stability and requiring multi-stage training procedures.
Conversion with fine-tuning~\citep{rathi2020hybrid} and direct training approaches
such as SpikeLM~\citep{xing2024spikelm} rely on soft substitutes to maintain
differentiability, but inevitably introduce approximation errors that diffuse across
depth. The shared limitation is that surrogate gradients are applied globally,
without any guarantee that the forward and backward paths are mutually consistent.

% \noindent\textbf{The shared limitation.}
% Across all three axes, the underlying assumption is the same: approximation is
% unavoidable, so the goal is to minimize it. This assumption is precisely what LASER
% challenges. By establishing an exact bijective encoding for linear computation (BSE)
% and confining nonlinear approximation to a single bounded module (ASNC), the forward
% path becomes provably consistent with the original ANN, transforming STE from a
% heuristic into a principled procedure. The distinction is not one of degree but of
% kind: prior methods approximate everywhere and hope errors stay small; LASER
% eliminates approximation where it is avoidable and bounds it where it is not.


%%%%%%%%%%%%%%%%%%%%%%%%%%%%%%%%%%%%%%%%%%%%%%%%%%%%%%%%%%%%%%%%%%%%%%%%%%%
\section{Preliminaries}
\label{sec:preliminaries}
%%%%%%%%%%%%%%%%%%%%%%%%%%%%%%%%%%%%%%%%%%%%%%%%%%%%%%%%%%%%%%%%%%%%%%%%%%%

Before presenting LASER, we establish the computational model, formalize the
conversion problem, and derive the error accumulation bound that governs why existing
methods fail at scale. Each element introduced here is used directly in the
methodology that follows.

\noindent\textbf{SNN Computation Model.}
\label{subsec:snn_model}
A spiking neural network communicates between layers via binary spike trains rather
than continuous activations. We adopt the non-leaky Integrate-and-Fire (IF) model
with soft reset and dynamic threshold as the computational primitive throughout our
framework. Given input current $I_{\mathrm{in}}(t)$, membrane potential $V_m(t)$,
and dynamic threshold $\Theta(t)$, the neuron emits spike $S(t) \in \{0,1\}$ and
updates its state according to
\begin{equation}
\label{eq:if_forward}
S(t) = \mathbb{1}\!\left[V_m(t) + I_{\mathrm{in}}(t) \ge \Theta(t)\right], \qquad
V_m(t{+}1) = V_m(t) + I_{\mathrm{in}}(t) - S(t)\,\Theta(t),
\end{equation}
initialized with $V_m(0) = 0$. The dynamic threshold $\Theta(t)$ is the key degree
of freedom: by setting it to a fixed positional schedule, the same IF unit realizes
exact bit-serial decomposition (BSE); by making it learnable, it realizes piecewise
nonlinear approximation (ASNC). A layer's full output is the spike train
$\mathbf{S} = (S(0), \dots, S(T-1)) \in \{0,1\}^T$.

\noindent\textbf{Formal Problem Statement.}
\label{subsec:problem}
Let $\mathcal{F}_{\mathrm{ANN}}$ be a pretrained ANN with $L$ layers, weights
$\{\mathbf{W}^\ell\}_{\ell=1}^L$, and intermediate activations
$\mathbf{x}^\ell \in \mathbb{R}^{d_\ell}$. The network alternates between affine
operations $\mathbf{z}^\ell = \mathbf{W}^\ell \mathbf{x}^{\ell-1} + \mathbf{b}^\ell$
and nonlinear activations $\mathbf{x}^\ell = \sigma(\mathbf{z}^\ell)$, producing
output $\mathbf{y} = \mathcal{F}_{\mathrm{ANN}}(\mathbf{x}^0)$.
%
ANN-to-SNN conversion seeks a spiking counterpart $\mathcal{F}_{\mathrm{SNN}}$ such
that
\begin{equation}
\label{eq:conversion_goal}
\bigl\|\mathcal{F}_{\mathrm{SNN}}(\mathbf{x}^0)
      - \mathcal{F}_{\mathrm{ANN}}(\mathbf{x}^0)\bigr\| \;\le\; \epsilon_{\mathrm{tol}},
\end{equation}
subject to the constraint that all inter-layer communication in
$\mathcal{F}_{\mathrm{SNN}}$ uses binary spike trains $\mathbf{S} \in \{0,1\}^T$.
The goal is to make $\epsilon_{\mathrm{tol}}$ as small as possible without
sacrificing the computational advantages of spiking hardware. We will show that
LASER achieves $\epsilon_{\mathrm{tol}}$ bounded by a quantity independent of depth
$L$, whereas all prior methods incur a bound that grows exponentially with $L$.

\noindent\textbf{Error Accumulation.}
\label{subsec:error_accumulation}
To understand why existing conversion methods fail on large models, consider what
happens when each layer introduces a relative approximation error. Let
$\hat{\mathbf{z}}^\ell$ denote the spike-domain approximation of the true linear
output $\mathbf{z}^\ell$, with per-layer error
$\|\hat{\mathbf{z}}^\ell - \mathbf{z}^\ell\| \le \epsilon_\ell
\|\mathbf{z}^\ell\|$.
After passing through the nonlinear activation $\sigma$ with Lipschitz constant
$L_\sigma$, the error at the output of layer $\ell$ satisfies
\[
\|\hat{\mathbf{x}}^\ell - \mathbf{x}^\ell\|
\;\le\; L_\sigma \epsilon_\ell \|\mathbf{z}^\ell\|.
\]
This error becomes the input perturbation for the next layer. Propagating through
all $L$ layers, the total output error grows as
\begin{equation}
\label{eq:error_accumulation}
\bigl\|\hat{\mathbf{y}} - \mathbf{y}\bigr\|
\;\le\; \|\mathbf{y}\| \cdot \prod_{\ell=1}^{L}(1 + L_\sigma \epsilon_\ell)
\;=\; O\!\left((1 + \epsilon)^L\right),
\end{equation}
where $\epsilon = \max_\ell \epsilon_\ell$. This is an exponential bound in depth
$L$. For a large language model with hundreds of transformer layers, even a modest
per-layer error $\epsilon$ produces catastrophic total degradation. The bound makes
precise what we observed empirically in the introduction: prior SNN methods are
structurally incapable of converting deep networks accurately.

\begin{remark}
\label{rem:error_structure}
Equation~\eqref{eq:error_accumulation} reveals the precise path to accurate
large-scale conversion. The exponential growth is driven by two factors: the number
of layers $L$ and the per-layer error $\epsilon_\ell$. Reducing $\epsilon_\ell$
uniformly across all layers requires prohibitive compute and still leaves exponential
scaling. A more effective strategy is to identify which layers can achieve
$\epsilon_\ell = 0$ exactly, and bound $\epsilon_\ell$ tightly only where it is
unavoidable. This is the design principle of LASER: BSE makes $\epsilon_\ell = 0$
for all linear layers, and ASNC bounds $\epsilon_\ell$ for nonlinear layers. The
product in~\eqref{eq:error_accumulation} then collapses from exponential in $L$
to linear in $M$, the number of nonlinear modules.
\end{remark}

This remark is not merely motivational — it is the formal blueprint for the
methodology. Section~\ref{subsec:bse} establishes the $\epsilon_\ell = 0$ claim for
linear layers, Section~\ref{subsec:asnc} establishes the bound for nonlinear layers,
and Section~\ref{subsec:ste} shows the gradient follows the same decomposition.


%%%%%%%%%%%%%%%%%%%%%%%%%%%%%%%%%%%%%%%%%%%%%%%%%%%%%%%%%%%%%%%%%%%%%%%%%%%
\section{Methodology}
\label{sec:methodology}
%%%%%%%%%%%%%%%%%%%%%%%%%%%%%%%%%%%%%%%%%%%%%%%%%%%%%%%%%%%%%%%%%%%%%%%%%%%

LASER converts a pretrained ANN to an SNN by systematically addressing each term in
the error accumulation bound~\eqref{eq:error_accumulation}. The three components
are not independent engineering choices but a coherent chain: each one is the logical
prerequisite for the next, and together they replace the exponential bound with one
that is linear in $M$ and independent of depth $L$.

\subsection{Bit Spike Encoding (BSE): Eliminating Linear Error}
\label{subsec:bse}

The first step is to make $\epsilon_\ell = 0$ for all linear layers in
Eq.~\eqref{eq:error_accumulation}. Existing encoding schemes map continuous
activations to spike trains through approximation — rate coding accumulates spike
counts, TTFS measures spike timing — and both incur irreducible encoding error that
enters the product in~\eqref{eq:error_accumulation} at every layer. The question is
whether it is possible to represent a numerical value as a spike train with zero
information loss. The answer is yes, and the construction is simpler than one might
expect.

An $N$-bit integer $q \in \{0, \dots, 2^N-1\}$ has a unique binary representation
$q = \sum_{t=0}^{N-1} b_t\,2^{N-1-t}$. We can unfold this representation across
time: assign timestep $t$ to bit $b_t$, emit a spike at timestep $t$ if and only if
$b_t = 1$, and recover $q$ exactly from the spike train by weighted summation. This
is the essence of BSE — it reframes spike encoding not as approximation but as
bit-serial communication.
%
Formally, we realize this with the IF model~\eqref{eq:if_forward} under the
positional threshold schedule
\begin{equation}
\label{eq:bse_schedule}
\Theta(t) = W_{\mathrm{bit}}(t) = 2^{N-1-t}, \qquad t = 0, \dots, N-1,
\end{equation}
driven by a single input impulse $I_{\mathrm{in}}(t) = \delta_{t0}\,q$ with
$V_m(0) = 0$. The BSE encoder $\mathsf{E}_N : \mathcal{Q}_N \to \{0,1\}^N$ is the
resulting spike train, and the BSE decoder $\mathsf{D}_N : \{0,1\}^N \to
\mathcal{Q}_N$ recovers $q$ via
\begin{equation}
\label{eq:bse_decoder}
\mathsf{D}_N\!\big(\{S(t)\}\big) = \sum_{t=0}^{N-1} S(t)\,2^{N-1-t}.
\end{equation}
The specific choice of threshold schedule~\eqref{eq:bse_schedule} is what makes the
IF model emit bit $b_t$ at exactly timestep $t$, as the following lemma establishes.

\begin{lemma}[Bit-exact emission]
\label{lem:bit_exact_main}
Under the threshold schedule~\eqref{eq:bse_schedule} and impulse input
$I_{\mathrm{in}}(t) = \delta_{t0}\,q$, the IF model~\eqref{eq:if_forward} emits
$S(t) = b_t$ for all $t = 0, \dots, N-1$. The membrane potential satisfies
$V_m(t) = \sum_{\tau > t} b_\tau\,2^{N-1-\tau}$, tracking the value of the
remaining unprocessed bits at each step.
\end{lemma}

The proof proceeds by induction on $t$ and is given in full in
Appendix~\ref{app:bse_entropy_proof}. The membrane potential invariant is the key
insight: at each step, $V_m(t)$ holds exactly the residual value after stripping off
the most significant bits already emitted, so the threshold comparison recovers the
next bit deterministically.

With Lemma~\ref{lem:bit_exact_main} in hand, we can establish the two properties
that make BSE the foundation of LASER.

\begin{theorem}[BSE bijectivity and entropy preservation]
\label{thm:bse_main}
The encoder $\mathsf{E}_N$ and decoder $\mathsf{D}_N$ are mutual inverses:
$\mathsf{D}_N(\mathsf{E}_N(q)) = q$ for all $q \in \mathcal{Q}_N$ and
$\mathsf{E}_N(\mathsf{D}_N(\mathbf{S})) = \mathbf{S}$ for all
$\mathbf{S} \in \{0,1\}^N$. Consequently, $H(\mathsf{E}_N(Q)) = H(Q)$ for any
integer-valued random variable $Q$: BSE preserves Shannon entropy exactly.
\end{theorem}

Bijectivity follows directly from Lemma~\ref{lem:bit_exact_main}; entropy
preservation follows because a bijection between finite sets is a relabeling that
leaves the probability distribution structurally unchanged. Full proof in
Appendix~\ref{app:bse_entropy_proof}. The entropy preservation result confirms that
no information is lost in transit: the spike train carries exactly as much
information as the original integer, no more and no less.

Theorem~\ref{thm:bse_main} establishes that BSE is lossless. The following corollary
connects this directly to the error accumulation bound~\eqref{eq:error_accumulation}.

\begin{corollary}[Linear layers contribute zero error]
\label{cor:no_diffusion_main}
Under $N$-bit alignment, the BSE forward pass through any linear layer is
functionally identical to the corresponding $N$-bit quantized ANN forward pass. Any
discrepancy from full-precision arithmetic equals the ANN's own quantization error
and does not diffuse or amplify. In terms of Eq.~\eqref{eq:error_accumulation},
$\epsilon_\ell = 0$ for every linear layer.
\end{corollary}

The proof uses the fact that the Soma accumulation $Q_a = \sum_t S_a(t)\,
W_{\mathrm{bit}}(t)$ is identical to BSE decoding, so weight-activation
multiplication is exact integer arithmetic. Full derivation and the Soma construction
are given in Appendix~\ref{app:bse_entropy_proof}. Setting $\epsilon_\ell = 0$ for
all linear layers in~\eqref{eq:error_accumulation} eliminates those terms from the
product entirely. What remains is the contribution of nonlinear layers — which is
where ASNC takes over.

\noindent\textbf{Activation-activation products.} For operations involving two spike-encoded activations, such as the Query-Key product in self-attention, direct BSE-domain computation is not applicable. We handle these via a decode-compute-re-encode strategy: both spike trains are losslessly decoded, the product is computed with standard multiply-accumulate operations, and the result is immediately re-encoded into BSE format. This preserves mathematical equivalence to the quantized ANN and ensures the output is a valid BSE spike train, so Corollary~\ref{cor:no_diffusion_main} applies to all downstream layers.

\subsection{Adaptive Spiking Neural Codec (ASNC): Bounding Nonlinear Error}
\label{subsec:asnc}

Corollary~\ref{cor:no_diffusion_main} eliminates $\epsilon_\ell$ for linear layers,
but nonlinear activation functions such as SiLU and Softmax have no closed-form
spike-domain equivalent. Some approximation at these points is unavoidable, and the
question becomes: how tightly can we bound $\epsilon_\ell$ for nonlinear layers, and
can we prevent that error from re-entering the linear layers that follow? ASNC answers
both questions.

The key design constraint is that ASNC must produce outputs in BSE format. If it
does, then Corollary~\ref{cor:no_diffusion_main} guarantees that subsequent linear
layers process those outputs exactly, and the nonlinear approximation error is frozen
at the point of re-encoding. This is the localization mechanism: BSE acts as a
firewall between ASNC's approximation error and the rest of the network.

Rather than replacing the nonlinearity with a spike-friendly surrogate, ASNC
approximates the target function $\sigma$ directly in the spike domain via
distribution-aware piecewise adaptation. The input domain is partitioned into $K$
segments. Since activation inputs in neural networks follow a bounded near-Gaussian
distribution~\citep{he2024understanding,lin2024duquant}, segments are allocated
non-uniformly: high-density and outlier regions receive finer partitions for
higher-fidelity fitting, while low-density regions receive coarser partitions to
reduce overhead. For segment $i$ covering interval $[a_i, b_i]$, a compact spiking
codec maps the input to a decoded analog estimate using the same IF
model~\eqref{eq:if_forward}, now with learned parameters:
\begin{equation}
\label{eq:asnc_codec}
\mathrm{spikes}_i(t)
  = \mathrm{LIF}\!\bigl(\mathrm{quantize}(x\,s_i + c_i),\,\theta_i(t)\bigr),
\qquad
\mathcal{C}_i(x) = \sum_{t=1}^{T} w_i(t)\,\mathrm{spikes}_i(t),
\end{equation}
where $s_i, c_i$ are affine parameters, $\theta_i(t)$ is the learned firing
threshold, and $w_i(t)$ are trainable decoding weights. The output is immediately
re-encoded into BSE format:
\begin{equation}
\label{eq:asnc_output}
\mathbf{S}^{\mathrm{BSE}}(t)
  = \sum_{i=1}^{K} \mathcal{I}_i(x)\,\mathbf{S}^{\mathrm{BSE}}_i(t),
\end{equation}
where $\mathcal{I}_i(x)$ is the segment selector. Note that
Eq.~\eqref{eq:asnc_output} produces a valid BSE spike train by construction, so
Corollary~\ref{cor:no_diffusion_main} applies immediately to all layers downstream of
ASNC. The approximation error introduced by ASNC is thus isolated at the nonlinear
module and cannot propagate further. Segments are split adaptively during training
when their normalized approximation difficulty exceeds a threshold; full details of
the splitting criterion and morphology-aware weighting are given in
Appendix~\ref{app:asnc_details}.

\noindent\textbf{Theoretical guarantees.}
Two results make the ASNC approximation precise. The first establishes that ASNC
achieves a strictly tighter error bound than prior uniform methods under the
activation distribution commonly observed in large networks.
 
\begin{theorem}[ASNC approximation to arbitrary precision]
\label{thm:asnc_bound}
For any target nonlinear function $\sigma$ and any $\delta > 0$, there exists a
segment resolution $K$ such that the ASNC approximation distortion satisfies
$D_K^{\mathrm{ASNC}} < \delta$. Moreover, under Gaussian activation inputs, the
optimal $K$-segment non-uniform partition achieves strictly lower asymptotic
distortion than any uniform $K$-segment scheme:
$D_K^{\mathrm{non-uniform}} < D_K^{\mathrm{uniform}}$ for all $K$.
\end{theorem}
The first claim follows from the universality of the LIF codec within each segment:
as $K$ grows, segments cover arbitrarily small intervals on which the codec converges
to the Lloyd-Max optimal reconstruction point. The second claim follows from
Bennett's integral and Holder's inequality applied to the Gaussian density. Full
proof in Appendix~\ref{app:asnc_proof}.
%
The second result connects ASNC's bounded error back to the accumulation
bound~\eqref{eq:error_accumulation}, completing the theoretical argument of LASER.
\begin{corollary}[ASNC error localization and linear depth bound]
\label{cor:asnc_local_main}
The ASNC approximation error $\epsilon_{\mathrm{ASNC}} \le D_K^{\mathrm{ASNC}}$
does not propagate beyond the nonlinear module. Because ASNC output is re-encoded
into BSE format via~\eqref{eq:asnc_output}, Corollary~\ref{cor:no_diffusion_main}
ensures all downstream linear layers operate exactly. The total conversion error
across a network with $M$ nonlinear modules satisfies
\[
\bigl\|\hat{\mathbf{y}} - \mathbf{y}\bigr\|
\;\le\; M \cdot \epsilon_{\mathrm{ASNC}},
\]
which is linear in $M$ and independent of network depth $L$.
\end{corollary}
 
This corollary directly answers the question raised by Remark~\ref{rem:error_structure}
in the Preliminaries. The exponential bound $O((1+\epsilon)^L)$
from~\eqref{eq:error_accumulation} is replaced with $O(M \cdot
\epsilon_{\mathrm{ASNC}})$: the terms for linear layers vanish (Corollary
\ref{cor:no_diffusion_main}), and the terms for nonlinear layers do not compound
(Corollary~\ref{cor:asnc_local_main}). For a transformer where $M \ll L$, this is an
exponential-to-linear improvement in how error scales with depth — the formal
explanation for why LASER maintains near-lossless accuracy at 70B scale where prior
methods fail catastrophically.

\subsection{STE Training: A Principled Gradient Procedure}
\label{subsec:ste}

With BSE and ASNC establishing high-fidelity forward propagation, the remaining
question is whether we can train ASNC's parameters stably. The non-differentiability
of the BSE encoder blocks standard backpropagation. In prior SNN work, this is
handled by surrogate gradients — smooth approximations of the non-differentiable spike
function — but these introduce gradient errors at every layer, compounding with depth
in exactly the same way as the forward errors of~\eqref{eq:error_accumulation}. The
key insight is that this problem has already been solved by the forward-path analysis:
since BSE makes linear layers exact and ASNC bounds nonlinear error, the gradient
need not be approximated globally. It can be handled with the same layer-type
decomposition.

ASNC parameters are adapted to a target pretrained model using a small calibration
set — we use MMLU — without modifying the original weights. For backpropagation
through the BSE encoder $S_{\mathrm{out}} = \mathsf{E}_N(V_{\mathrm{in}})$, we apply
the hard Straight-Through Estimator (STE), treating the encoder gradient as the
identity mapping $\partial S_{\mathrm{out}} / \partial V_{\mathrm{in}} = \mathbf{I}$.
%
In prior work, STE is a heuristic adopted because the forward pass is lossy and no
better alternative exists. In LASER, the forward-path fidelity guarantees of BSE and
ASNC give STE a precise theoretical justification.
\begin{proposition}[STE correctness under LASER]
\label{prop:ste_main}
Under the LASER framework:
\begin{enumerate}
  \item For all linear layers, the STE gradient equals the true gradient exactly —
        no approximation is introduced.
  \item For nonlinear layers, the STE gradient error is bounded by
        $\epsilon_{\mathrm{ASNC}} \le D_K^{\mathrm{ASNC}}$.
  \item The total gradient approximation error across the network is
        $O(M \cdot \epsilon_{\mathrm{ASNC}})$, independent of network depth $L$.
\end{enumerate}
\end{proposition}

The proof of (i) uses Corollary~\ref{cor:no_diffusion_main}: since the BSE forward
pass is the identity on the integer grid, the chain rule through a linear layer is
unaffected by the STE substitution. The proof of (ii) uses
Theorem~\ref{thm:asnc_bound}: the forward error at an ASNC module is bounded by
$D_K^{\mathrm{ASNC}}$, and the gradient error is proportional to this. The proof of
(iii) uses Corollary~\ref{cor:asnc_local_main}: since forward errors do not
accumulate across depth, neither do gradient errors. Full proof in
Appendix~\ref{app:ste_proof}.

\begin{remark}
Proposition~\ref{prop:ste_main} closes the loop on the conversion
problem~\eqref{eq:conversion_goal}. The forward pass satisfies
$\|\hat{\mathbf{y}} - \mathbf{y}\| \le M \cdot \epsilon_{\mathrm{ASNC}}$ by
Corollary~\ref{cor:asnc_local_main}, and the gradient procedure that minimizes
$\epsilon_{\mathrm{ASNC}}$ obeys the same depth-independent bound by
Proposition~\ref{prop:ste_main}. Both forward accuracy and training stability are
governed by a single controllable quantity — the ASNC segment resolution $K$ — which
can be increased to tighten the bound at the cost of additional ASNC capacity. The
conversion tolerance $\epsilon_{\mathrm{tol}}$ in~\eqref{eq:conversion_goal} is
therefore not a fixed property of the architecture but a tunable parameter under the
practitioner's control.
\end{remark}

\section{Experiments}
\label{sec:experiments}

Our experiments are organized to validate the three formal claims of
Section~\ref{sec:methodology} in sequence. Section~\ref{subsec:exp_bse} validates
Corollary~\ref{cor:no_diffusion_main}: BSE linear layers contribute $\epsilon_\ell =
0$. Section~\ref{subsec:exp_asnc} validates Theorem~\ref{thm:asnc_bound} and
Corollary~\ref{cor:asnc_local_main}: ASNC achieves arbitrary precision with increasing
$K$ and error does not propagate beyond the nonlinear module.
Section~\ref{subsec:exp_scale} validates the depth-independence of the bound: PPL
degradation does not grow with model depth. Section~\ref{subsec:exp_energy} reports
energy efficiency on neuromorphic hardware.
%
Unless otherwise stated, all experiments use LLaMA-2 7B on WikiText-2 with FP16
precision and a fixed 16-step budget. System-level perplexity is reported as mean
$\pm$ std over 5 seeds; the FP16 ANN baseline achieves PPL $= 5.12$.


%------------------------------------------------------------------
\subsection{Validating BSE: Linear Layers Contribute Zero Error}
\label{subsec:exp_bse}
%------------------------------------------------------------------

Corollary~\ref{cor:no_diffusion_main} predicts that BSE-encoded linear layers
introduce no approximation error beyond the ANN's own quantization floor. We test
this at two levels.

\noindent\textbf{Encoding fidelity.}
We measure reconstruction MSE for BSE, rate coding, and TTFS under matched timestep
budgets, using 10{,}000 random values per precision. As shown in
Figure~\ref{fig:bse_fidelity}, BSE achieves MSE of $1.03 \times 10^{-9}$ at FP16
and $1.33 \times 10^{-14}$ at FP32, matching machine precision. Rate coding and TTFS
under the same budgets yield MSE as large as $10^4$--$10^5$ — eleven to eighteen
orders of magnitude larger. Critically, TTFS requires $1{,}024$ steps to reach FP16
precision, a $64\times$ latency increase, whereas BSE reaches the same precision in
16 steps.

\noindent\textbf{Linear layer fidelity.}
We extract a representative FFN linear layer, pass 1{,}024 random inputs through the
FP16 ANN reference, and compare decoded SNN outputs via MSE. As shown in
Figure~\ref{fig:component_mse}(a), BSE-based linear operations achieve MSE of $1.45
\times 10^{-9}$, essentially identical to the INT16 quantization floor ($1.12 \times
10^{-9}$). Rate coding under the same setting degrades to $4.88 \times 10^{-1}$ —
eight orders of magnitude worse. This confirms that $\epsilon_\ell = 0$ for BSE
linear layers up to the quantization floor, as Corollary~\ref{cor:no_diffusion_main}
predicts.


%------------------------------------------------------------------
\subsection{Validating ASNC: Bounded and Localized Nonlinear Error}
\label{subsec:exp_asnc}
%------------------------------------------------------------------

Theorem~\ref{thm:asnc_bound} predicts that ASNC achieves arbitrary precision as $K$
increases, and Corollary~\ref{cor:asnc_local_main} predicts that this error does not
propagate beyond the nonlinear module.

\noindent\textbf{Arbitrary precision with increasing $K$.}
\textcolor{red}{\textbf{[TODO: Add figure showing ASNC approximation distortion
$D_K^{\mathrm{ASNC}}$ as a function of segment count $K$ for SiLU and Softmax, on a
log-log scale. Should show monotone decrease toward zero, validating
Theorem~\ref{thm:asnc_bound}. Also include a comparison curve for uniform
partitioning to empirically confirm $D_K^{\mathrm{non-uniform}} <
D_K^{\mathrm{uniform}}$.]}}

\noindent\textbf{Error localization.}
Figure~\ref{fig:component_mse}(a) shows that the only visible MSE in the full FFN
block originates from ASNC, bounded at $8.73 \times 10^{-7}$. The BSE linear path
contributes nothing beyond the quantization floor. To confirm that ASNC error does
not propagate, we measure PPL after replacing individual components with low-fidelity
alternatives. As shown in Figure~\ref{fig:component_mse}(b), replacing ASNC with an
inconsistent ReLU degrades PPL catastrophically to $50.74$ and replacing BSE with
rate coding leads to failure beyond $100$ PPL, while our BSE+ASNC adds only $+0.09$
over the ideal ceiling of BSE with the original SiLU. This confirms that both
components are load-bearing: BSE provides the lossless linear path, and ASNC provides
the localized nonlinear approximation.

\begin{figure}[h!]
\centering
\begin{minipage}[c]{0.46\textwidth}
\centering
\begin{tikzpicture}
\begin{axis}[
    width=\textwidth,
    height=\textwidth,
    xlabel={\scriptsize MSE (log scale)},
    xmode=log,
    xmin=5e-12, xmax=5e2,
    ymin=0.3, ymax=5.7,
    ytick={1,2,3,4,5},
    yticklabels={
        {\scriptsize INT16 Quant.},
        {\scriptsize BSE Linear},
        {\scriptsize ASNC SiLU},
        {\scriptsize Full FFN},
        {\scriptsize Rate Coding}
    },
    y dir=reverse,
    tick label style={font=\scriptsize},
    label style={font=\scriptsize},
    axis line style={thin, black!70},
    major tick length=1.5pt,
    xmajorgrids=true,
    grid style={TolGrey!25, thin},
    axis on top,
    clip=false,
    axis background/.style={fill=black!4},
]
\fill[TolGrey!50, draw=TolGrey!70]
    (axis cs:5e-12, 0.72) rectangle (axis cs:1.12e-9, 1.28);
\fill[TolBlue!60, draw=TolBlue]
    (axis cs:5e-12, 1.72) rectangle (axis cs:1.45e-9, 2.28);
\fill[TolBlue!60, draw=TolBlue]
    (axis cs:5e-12, 2.72) rectangle (axis cs:8.73e-7, 3.28);
\fill[TolBlue!60, draw=TolBlue]
    (axis cs:5e-12, 3.72) rectangle (axis cs:9.51e-7, 4.28);
\fill[TolRed!60, draw=TolRed]
    (axis cs:5e-12, 4.72) rectangle (axis cs:4.88e-1, 5.28);
\node[font=\tiny, anchor=west] at (axis cs:3e-9,  1) {$10^{-9}$};
\node[font=\tiny, anchor=west] at (axis cs:3e-9,  2) {$10^{-9}$};
\node[font=\tiny, anchor=west] at (axis cs:2e-6,  3) {$10^{-7}$};
\node[font=\tiny, anchor=west] at (axis cs:2e-6,  4) {$10^{-7}$};
\node[font=\tiny, anchor=west] at (axis cs:1.2e0, 5) {$10^{-1}$};
\end{axis}
\end{tikzpicture}
\end{minipage}
\hfill
\begin{minipage}[c]{0.50\textwidth}
\centering
\small
\begin{tabular}{lc}
\toprule
\textbf{FFN Configuration} & \textbf{PPL} \\
\midrule
ANN with SiLU (Baseline)       & $5.12 \pm 0.01$ \\
\textit{BSE + Standard SiLU}   & \textit{$5.18 \pm 0.01$} \\
\textbf{BSE + ASNC (Ours)}     & $\mathbf{5.21 \pm 0.02}$ \\
BSE + ReLU                     & $50.74 \pm 0.04$ \\
Rate Coding + ASNC             & $103.5 \pm 10.4$ \\
\bottomrule
\end{tabular}
\end{minipage}
\caption{\textbf{(a)} Component-level MSE (FP16, 16 steps). BSE linear layers match
the INT16 quantization floor; ASNC confines nonlinear error to $10^{-7}$; rate coding
degrades to $10^{-1}$. \textbf{(b)} System-level PPL on LLaMA-2 7B. BSE+ASNC is
only $+0.09$ above the ideal ceiling; removing either component causes catastrophic
degradation, confirming both are necessary.}
\label{fig:component_mse}
\end{figure}

\noindent\textbf{Ablation: layer-wise and progressive SNN-ization.}
To further confirm error localization, we SNN-ize one subsystem at a time and measure
the resulting PPL. Table~\ref{tab:ablation} (left) shows that the dominant
sensitivity comes from FFN layers ($+0.23$) and attention layers ($+0.19$), both of
which contain nonlinear operations, while the nearly-linear LayerNorm and Embedding
contribute only $+0.06$ and $+0.08$. Table~\ref{tab:ablation} (right) shows the
progressive composition: errors accumulate additively and saturate well below one
point above baseline, consistent with the $(1 + L_\sigma\delta)^M$ bound of
Corollary~\ref{cor:asnc_local_main} — errors do not compound across layers.

\begin{table}[h!]
\centering
\caption{Ablation on WikiText-2 PPL (FP16, 16 steps; ANN baseline $= 5.12 \pm
0.01$). \textbf{Left:} Layer-wise sensitivity shows nonlinear modules dominate.
\textbf{Right:} Progressive SNN-ization shows additive, bounded accumulation
consistent with Corollary~\ref{cor:asnc_local_main}.}
\label{tab:ablation}
\vspace{0.4em}
\begin{minipage}[t]{0.48\textwidth}
\centering
\small
\begin{tabular}{lc}
\toprule
\textbf{SNN-ized Component} & \textbf{PPL} \\
\midrule
None (Baseline)           & $5.12 \pm 0.01$ \\
FFN Layers only           & $5.35 \pm 0.03$ \\
Attention Layers only     & $5.31 \pm 0.02$ \\
Embedding/Output only     & $5.20 \pm 0.02$ \\
LayerNorm only            & $5.18 \pm 0.01$ \\
Full SNN (All)            & $5.58 \pm 0.04$ \\
\bottomrule
\end{tabular}
\end{minipage}
\hfill
\begin{minipage}[t]{0.48\textwidth}
\centering
\small
\begin{tabular}{lc}
\toprule
\textbf{Progressive Strategy} & \textbf{PPL} \\
\midrule
FFN only                & $5.35 \pm 0.03$ \\
$+$ Activation          & $5.47 \pm 0.04$ \\
$+$ Attention           & $5.50 \pm 0.04$ \\
$+$ Embedding           & $5.58 \pm 0.04$ \\
Full SNN                & $5.61 \pm 0.04$ \\
\bottomrule
\end{tabular}
\end{minipage}
\end{table}


%------------------------------------------------------------------
\subsection{Validating Depth-Independence: Scalability Across Model Sizes}
\label{subsec:exp_scale}
%------------------------------------------------------------------

Corollary~\ref{cor:asnc_local_main} predicts that the conversion error bound
$(1 + L_\sigma\delta)^M$ is independent of network depth $L$. A direct empirical
consequence is that PPL degradation should not increase — and may in fact decrease —
as models grow deeper, since larger models tend to be more robust to bounded
perturbations.

\noindent\textbf{Depth-independence across scales.}
\textcolor{red}{\textbf{[TODO: Add a figure plotting PPL degradation ($\Delta$PPL =
SNN PPL $-$ ANN PPL) on the y-axis against model depth (number of transformer layers)
on the x-axis, for all five models in Table~\ref{tab:scaling}. This directly
visualizes whether $\Delta$PPL grows with depth (prior methods) or remains flat /
decreases (LASER). Expected result: flat or decreasing trend, consistent with
depth-independent bound.]}}

\noindent\textbf{End-to-end accuracy across model families.}
Table~\ref{tab:scaling} reports task accuracy across five LLMs under a fixed 16-step
budget. SNN performance consistently tracks ANN baselines across all benchmarks and
model sizes. On LLaMA-2 70B, the largest model tested, MMLU accuracy drops by only
$1.2\%$ and overall degradation across four benchmarks remains below $2\%$. Notably,
degradation on the 70B model is comparable to or smaller than on 7B, consistent with
the depth-independence prediction.

\begin{table}[h!]
\centering
\caption{End-to-end performance (accuracy, \%; FP16, 16-step budget). Degradation
from ANN to SNN remains within $2\%$ across all models and benchmarks. Degradation
does not grow with model scale, consistent with the depth-independent bound of
Corollary~\ref{cor:asnc_local_main}.}
\label{tab:scaling}
\small
\begin{tabular*}{\textwidth}{@{\extracolsep{\fill}}l cc cc cc cc}
\toprule
\multirow{2}{*}{Model}
  & \multicolumn{2}{c}{MMLU}
  & \multicolumn{2}{c}{HellaSwag}
  & \multicolumn{2}{c}{ARC}
  & \multicolumn{2}{c}{TruthfulQA} \\
\cmidrule(lr){2-3}\cmidrule(lr){4-5}\cmidrule(lr){6-7}\cmidrule(lr){8-9}
  & ANN & SNN & ANN & SNN & ANN & SNN & ANN & SNN \\
\midrule
Phi-2 (2.7B)          & 58.11 & 56.0 & 75.11 & 72.5 & 61.09 & 58.9 & 44.47 & 42.1 \\
LLaMA-2 (7B)          & 60.04 & 58.2 & 79.13 & 76.9 & 56.14 & 54.5 & 40.95 & 39.0 \\
Mistral (7.3B)        & 60.78 & 59.1 & 84.88 & 83.0 & 63.14 & 61.5 & 68.26 & 66.0 \\
Mistral (8$\times$7B) & 68.59 & 68.0 & 86.03 & 85.2 & 67.24 & 66.5 & 59.54 & 58.3 \\
LLaMA-2 (70B)         & 65.40 & 64.2 & 86.90 & 85.5 & 67.20 & 65.8 & 44.90 & 43.1 \\
\bottomrule
\end{tabular*}
\end{table}

\noindent\textbf{Comparison with prior SNN methods.}
Table~\ref{tab:comparison} compares LASER against SpikeLLM~\citep{xing2025spikellm}
on LLaMA-2 7B. Note that SpikeLLM operates under a different and harder regime —
W4A4 and W2A16 weight quantization — whereas LASER uses FP16 weights throughout.
Even under this conservative comparison, our PPL degradation of $+0.46$ is
substantially smaller than SpikeLLM's $+6.7$ to $+9.0$.

\textcolor{red}{\textbf{[TODO: If a same-precision (FP16 weights) baseline from prior
SNN work can be identified, add it here for a fully controlled comparison. If not,
keep the current comparison and the regime clarification.]}}

\begin{table}[h!]
\centering
\caption{Comparison on LLaMA-2 7B, WikiText-2. LASER uses FP16 weights; SpikeLLM
uses aggressive weight quantization (W4A4, W2A16), a harder and different setting.
Even under this conservative comparison, LASER degrades by $+0.46$ vs.\ $+6.7$ to
$+9.0$ for SpikeLLM.}
\label{tab:comparison}
\small
\begin{tabular}{lccc}
\toprule
\textbf{Method} & \textbf{Weight Precision} & \textbf{Steps} & \textbf{PPL} \\
\midrule
ANN Baseline              & FP16  & ---  & $5.12$ \\
SpikeLLM                  & W2A16 & 2--4 & $14.16$ \\
SpikeLLM                  & W4A4  & 2--4 & $11.85$ \\
\textbf{LASER (Ours)}     & FP16  & 16   & $\mathbf{5.58 \pm 0.04}$ \\
\bottomrule
\end{tabular}
\end{table}

\noindent\textbf{Fine-grained attention analysis.}
\textcolor{red}{\textbf{[TODO: Decide whether the attention fine-grained table
(Q/K/V vs.\ score path sensitivity) is needed in the main paper. If yes, move here as
supporting evidence for the linear-layer zero-error claim: Q/K/V projections are
linear and should contribute small, bounded error; the score path involves the QK
product which uses decode-compute-re-encode and may contribute more. This would
directly validate the different treatment of weight-activation vs.\
activation-activation products described in Section~\ref{subsec:bse}.]}}


%------------------------------------------------------------------
\subsection{Energy Efficiency on Neuromorphic Hardware}
\label{subsec:exp_energy}
%------------------------------------------------------------------

We benchmark the ASNC nonlinear component on Intel Loihi~2 against an ANN SiLU
baseline on an NVIDIA A100 under matched precisions. Loihi~2 energy per operation is
obtained from on-board probes; GPU energy is estimated from sustained-load power
measurements over millions of identical operations via \texttt{nvidia-smi}. As shown
in Table~\ref{tab:energy}, Loihi~2 consumes only ${\sim}0.5\%$ of the GPU energy at
both 16-bit and 32-bit settings, a more than $200\times$ reduction. This measurement
is for the nonlinear component specifically; system-level energy on full neuromorphic
deployment would depend on hardware availability for all layer types.

\textcolor{red}{\textbf{[TODO: If full-model Loihi~2 energy measurements become
available, add a row or separate table showing system-level energy ratio. This would
strengthen the hardware efficiency story significantly.]}}

\begin{table}[h!]
\centering
\caption{Energy ratio of a single ASNC nonlinear operation on Intel Loihi~2 vs.\
SiLU on NVIDIA A100. Loihi~2 requires only ${\sim}0.5\%$ of GPU energy, a more than
$200\times$ reduction at the component level.}
\label{tab:energy}
\small
\begin{tabular}{lcc}
\toprule
\textbf{Precision} & \textbf{Steps} & \textbf{Energy Ratio (Loihi~2 / GPU)} \\
\midrule
16-bit & 16 & $5.5 \times 10^{-3}$ \\
32-bit & 32 & $5.3 \times 10^{-3}$ \\
\bottomrule
\end{tabular}
\end{table}

\section{Conclusion}
\label{sec:conclusion}
 
LASER demonstrates that the long-standing accuracy gap between ANNs and SNNs at scale
is not an inherent consequence of spike-based computation, but a consequence of
treating conversion as an approximation problem rather than a representation problem.
By establishing an exact bijective encoding for linear computation and a provably
bounded, localized approximation for nonlinear computation, the exponential error
accumulation that has made large-scale SNN deployment impractical collapses to a
linear bound independent of network depth. The implication is broader than a single
framework: it suggests that the field's reliance on surrogate gradients and
spike-friendly architectural substitutes has been addressing symptoms rather than the
underlying cause, and that principled representation design is a more productive path
forward. Whether this principle extends to other neural coding schemes, other hardware
substrates, or online learning settings remains an open and promising direction.

\bibliographystyle{plainnat}
\bibliography{references}

%%%%%%%%%%%%%%%%%%%%%%%%%%%%%%%%%%%%%%%%%%%%%%%%%%%%%%%%%%%%%%%%%%%%%%%%%%%%%%%
% APPENDIX
%%%%%%%%%%%%%%%%%%%%%%%%%%%%%%%%%%%%%%%%%%%%%%%%%%%%%%%%%%%%%%%%%%%%%%%%%%%%%%%
\appendix

%%%%%%%%%%%%%%%%%%%%%%%%%%%%%%%%%%%%%%%%%%%%%%%%%%%%%%%%%%%%%%%%%%%%%%%%%%%
\section{BSE: Full Construction and Proofs}
\label{app:bse_entropy_proof}
%%%%%%%%%%%%%%%%%%%%%%%%%%%%%%%%%%%%%%%%%%%%%%%%%%%%%%%%%%%%%%%%%%%%%%%%%%%

This appendix provides the complete construction of the BSE encoder and decoder,
including the Soma linear layer mechanism referenced in Section~\ref{subsec:bse},
and the full proofs of Lemma~\ref{lem:bit_exact_main},
Theorem~\ref{thm:bse_main}, and Corollary~\ref{cor:no_diffusion_main}.

\subsection{Soma: Linear Layer Computation in the BSE Domain}
\label{app:soma}

The main paper describes BSE encoding and decoding at the level of a single integer
$q$. In a full network layer, weights and biases must also be absorbed into the
spike-domain computation. The \textbf{Soma} compartment realizes this: it accumulates
an incoming BSE spike train, applies a scalar weight and bias, re-normalizes the
result, and emits a new BSE-encoded spike train. This makes weight-activation
multiplication a native spike-domain operation requiring no decode-compute-re-encode
detour.

For a scalar weight $w \in \mathbb{R}$, bias $b \in \mathbb{R}$, input scale
$\lambda \in \mathbb{R}_{>0}$, and input offset $\mu \in \mathbb{R}$, the Soma
forward pass on an incoming BSE spike train $\{S_a(t)\}_{t=0}^{N-1}$ is:
\begin{equation}
\label{eq:soma_def_app}
\begin{aligned}
Q_a &= \sum_{t=0}^{N-1} S_a(t)\,W_{\mathrm{bit}}(t), \\
Y   &= \frac{Q_a}{\lambda}\,w - w\mu + b, \\
\mu_{\mathrm{out}} &= -\min_{y \in \mathcal{S}}\,y,
\qquad
R_{\mathrm{out}} = \max_{y \in \mathcal{S}}\,(y + \mu_{\mathrm{out}}) + \varepsilon, \\
\lambda_{\mathrm{out}} &= \frac{2^N - 1}{R_{\mathrm{out}}}, \\
Q_y &= \operatorname{clip}_{[0,\,2^N-1]}\!\left(
         \mathrm{round}\!\left(\lambda_{\mathrm{out}}\,(Y + \mu_{\mathrm{out}})\right)
       \right),
\end{aligned}
\end{equation}
where $W_{\mathrm{bit}}(t) = 2^{N-1-t}$ is the bit weight at step $t$, $\mathcal{S}$
is the set of observed $Y$ values used for range estimation (computed from a
calibration pass), and $\varepsilon > 0$ is a small stability constant. The output
integer $Q_y$ is then passed to the BSE encoder~\eqref{eq:standard_bit_schedule} to
produce the output spike train $\{S_y(t)\}_{t=0}^{N-1}$.

The key observation is that $Q_a = \mathsf{D}_N(\{S_a(t)\})$ by the definition of
the BSE decoder~\eqref{eq:bse_decoder}. Therefore the Soma accumulation step is
identical to lossless decoding, and the weight multiplication $Y$ operates on the
exact integer representation of the input. This establishes the exactness of
weight-activation multiplication in the BSE domain, which underlies the
Exactness of Spike-domain Linear Accumulation paragraph below.

\subsection{Setting and Notation}
\label{app:bse_setting}

Fix a bit-window length $N \in \mathbb{N}$. For each channel, let the per-layer
affine calibration be $(\lambda, \mu)$ with $\lambda > 0$ and $\mu \in \mathbb{R}$.
Define the $N$-bit quantizer
\begin{equation}
\label{eq:quantizer_grid_def}
\mathcal{Q}_N \;\triangleq\; \{0, 1, \dots, 2^N - 1\},
\qquad
q \;=\; \mathrm{clip}_{[0,\,2^N-1]}\!\Big(\mathrm{round}\big(\lambda\,(x+\mu)\big)\Big).
\end{equation}
The corresponding de-quantization map is $\hat{x} = \lambda^{-1}q - \mu$. Write the
binary expansion of $q$ as
\begin{equation}
\label{eq:binary_expansion_def}
q \;=\; \sum_{t=0}^{N-1} b_t\,2^{N-1-t}, \qquad b_t \in \{0,1\}.
\end{equation}

\subsection{BSE Encoder and Decoder}
\label{app:bse_encdec}

Let the bit-weights and thresholds follow the standard positional schedule
\begin{equation}
\label{eq:standard_bit_schedule}
W_{\mathrm{bit}}(t) \;=\; 2^{\,N-1-t},
\qquad
\Theta(t) \;=\; W_{\mathrm{bit}}(t),
\qquad t = 0, \dots, N-1.
\end{equation}
Drive the IF model~\eqref{eq:if_forward} with a single impulse:
\begin{equation}
\label{eq:scalar_drive}
I_{\mathrm{in}}(t) \;=\; \delta_{t0}\,q, \qquad V_m(0) = 0.
\end{equation}
The \emph{BSE encoder} $\mathsf{E}_N : \mathcal{Q}_N \to \{0,1\}^N$ is the spike
train $S(t)$ produced by~\eqref{eq:if_forward} for $t = 0, \dots, N-1$. The
\emph{BSE decoder} $\mathsf{D}_N : \{0,1\}^N \to \mathcal{Q}_N$ is the weighted sum
\begin{equation}
\label{eq:bse_decoder}
\mathsf{D}_N\!\big(\{s(t)\}_{t=0}^{N-1}\big)
\;=\; \sum_{t=0}^{N-1} s(t)\,2^{\,N-1-t}.
\end{equation}

\subsection{Proofs}
\label{app:bse_proofs}

\begin{lemma}[Bit-exact emission]
\label{lem:bit_exact}
Under~\eqref{eq:standard_bit_schedule} and~\eqref{eq:scalar_drive}, the encoder
emits the binary digits of $q$ in MSB-to-LSB order: $S(t) = b_t$ for all
$t = 0, \dots, N-1$, and the membrane satisfies
$V_m(t) = \sum_{\tau > t} b_\tau\,2^{\,N-1-\tau}$ for $t = 0, \dots, N$.
\end{lemma}

\begin{proof}
By~\eqref{eq:if_forward} and~\eqref{eq:scalar_drive}, $V_m(0) = 0$ and
$V_m(0) + I_{\mathrm{in}}(0) = q$. At $t = 0$ the threshold is
$\Theta(0) = 2^{N-1}$, so $M_0 = \mathbb{1}\{q \ge 2^{N-1}\} = b_0$ and the
post-update membrane is
\[
V_m(1) = q - b_0\,2^{N-1} = \sum_{\tau=1}^{N-1} b_\tau\,2^{\,N-1-\tau}.
\]
Assume the claim holds up to time $t$, so $V_m(t) = \sum_{\tau > t}
b_\tau\,2^{N-1-\tau}$ and $\Theta(t) = 2^{N-1-t}$. Then
\[
M_t
= \mathbb{1}\{V_m(t) \ge \Theta(t)\}
= \mathbb{1}\!\Big\{\sum_{\tau > t} b_\tau\,2^{N-1-\tau} \ge 2^{N-1-t}\Big\}
= b_t,
\]
since the leading term is $b_t\,2^{N-1-t}$ and all lower-order terms are strictly
smaller. The update gives
\[
V_m(t+1) = V_m(t) - M_t\,\Theta(t) = \sum_{\tau > t+1} b_\tau\,2^{N-1-\tau}.
\]
By induction the claim holds for all $t$, and $V_m(N) = 0$.
\end{proof}

This lemma corresponds to Lemma~\ref{lem:bit_exact_main} in the main paper.

\begin{theorem}[Encoder-decoder bijection]
\label{thm:bijective}
$\mathsf{E}_N$ and $\mathsf{D}_N$ are mutual inverses:
$\mathsf{D}_N(\mathsf{E}_N(q)) = q$ for all $q \in \mathcal{Q}_N$, and
$\mathsf{E}_N(\mathsf{D}_N(\mathbf{S})) = \mathbf{S}$ for all
$\mathbf{S} \in \{0,1\}^N$. In particular, $\mathsf{E}_N$ is a bijection.
\end{theorem}

\begin{proof}
The first identity follows from Lemma~\ref{lem:bit_exact} and~\eqref{eq:bse_decoder}:
\[
\mathsf{D}_N\!\big(\mathsf{E}_N(q)\big)
= \sum_{t=0}^{N-1} S(t)\,2^{N-1-t}
= \sum_{t=0}^{N-1} b_t\,2^{N-1-t} = q.
\]
For the second identity, any $\mathbf{S} \in \{0,1\}^N$ defines
$q^\star = \mathsf{D}_N(\mathbf{S})$ with binary digits $b_t^\star = S(t)$.
Lemma~\ref{lem:bit_exact} shows that driving the IF with $q^\star$ emits $b_t^\star$
at step $t$, i.e., $\mathsf{E}_N(q^\star) = \mathbf{S}$.
\end{proof}

\begin{theorem}[Entropy preservation under $N$-bit alignment]
\label{thm:entropy_preservation}
Let $X$ be any real-valued random variable and $Q = \mathrm{clip} \circ
\mathrm{round}(\lambda(X+\mu)) \in \mathcal{Q}_N$ its $N$-bit
quantization~\eqref{eq:quantizer_grid_def}. Let $\mathbf{S} = \mathsf{E}_N(Q)$.
Then $H(Q) = H(\mathbf{S})$.
\end{theorem}

\begin{proof}
By Theorem~\ref{thm:bijective}, $\mathsf{E}_N$ is a bijection between the finite sets
$\mathcal{Q}_N$ and $\{0,1\}^N$. Hence for all $s \in \{0,1\}^N$,
$\mathbb{P}[\mathbf{S} = s] = \mathbb{P}[Q = \mathsf{D}_N(s)]$. Therefore
\[
H(\mathbf{S})
= -\sum_{s} \mathbb{P}[\mathbf{S}=s] \log \mathbb{P}[\mathbf{S}=s]
= -\sum_{q} \mathbb{P}[Q=q] \log \mathbb{P}[Q=q]
= H(Q). \qedhere
\]
\end{proof}

Theorems~\ref{thm:bijective} and~\ref{thm:entropy_preservation} together constitute
the proof of Theorem~\ref{thm:bse_main} in the main paper.

\noindent\textbf{Vector and batch case.}
For $B$ samples and $d$ channels, apply $(\lambda, \mu)$ per channel and
$\mathsf{E}_N$ elementwise. The product map $\mathsf{E}_N^{\otimes Bd}$ is a
bijection between the two finite product sets, so the joint entropy is preserved:
$H(Q_{1:B,1:d}) = H(S_{0:N-1,1:B,1:d})$.

\noindent\textbf{Exactness of spike-domain linear accumulation.}
Within a bit window, the Soma pre-accumulation~\eqref{eq:soma_def_app} recovers the
decoded integer exactly:
\[
Q_a = \sum_{t=0}^{N-1} S_a(t)\,W_{\mathrm{bit}}(t) = \mathsf{D}_N\!\big(\{S_a(t)\}\big).
\]
For any weight $w$, the product $wQ_a$ is therefore identical to applying $w$ to the
decoded integer. All spike-domain linear accumulations are exact integer arithmetic.
The only approximation in the BSE pathway is the explicit rounding and clipping that
defines $Q$ or $Q_y$ in~\eqref{eq:soma_def_app}, which equals the ANN's own $N$-bit
quantization error.

\begin{corollary}[No error diffusion from BSE]
\label{cor:no_diffusion}
Assume each layer uses $(\lambda, \mu)$ matched to its $N$-bit quantization grid and
the binary schedule~\eqref{eq:standard_bit_schedule}. Then, relative to the
corresponding $N$-bit quantized ANN, the BSE forward pass is functionally identical:
encoding, spike-domain accumulation, and decoding are exact and introduce no
additional error. Any discrepancy from full-precision arithmetic equals the ANN's
own $N$-bit quantization error and neither diffuses nor amplifies through BSE.
\end{corollary}

This corollary is restated as Corollary~\ref{cor:no_diffusion_main} in the main
paper, where it is used to establish that $\epsilon_\ell = 0$ for all linear layers
in the error accumulation bound~\eqref{eq:error_accumulation}.

\noindent\textbf{Remarks.}
(i) The proof requires only a positional numeral system with unique representation;
the binary schedule~\eqref{eq:standard_bit_schedule} is the canonical choice.
(ii) Signed ranges are handled by the offset $\mu$, which translates the dynamic
range to $[0, 2^N-1]$ before encoding and back after decoding; bijectivity and
entropy preservation hold channelwise under this affine change of variables.
(iii) The results extend to per-channel $(\lambda_o, \mu_o)$ and per-layer bit
windows, provided $N$ is the operative grid width of the ANN baseline.


%%%%%%%%%%%%%%%%%%%%%%%%%%%%%%%%%%%%%%%%%%%%%%%%%%%%%%%%%%%%%%%%%%%%%%%%%%%
\section{ASNC: Error Bound and Localization Proofs}
\label{app:asnc_proof}
%%%%%%%%%%%%%%%%%%%%%%%%%%%%%%%%%%%%%%%%%%%%%%%%%%%%%%%%%%%%%%%%%%%%%%%%%%%

This appendix provides the full proofs of Theorem~\ref{thm:asnc_bound} and
Corollary~\ref{cor:asnc_local_main} in the main paper.

\subsection{Setup}

Let $X \sim \mathcal{N}(0,1)$ be the activation input. We consider the full real line
$(-\infty, \infty)$, reflecting that ASNC targets general nonlinearities such as SiLU
and Softmax defined on $\mathbb{R}$. The input domain is partitioned into $K$
intervals $\{I_i\}_{i=1}^K$, and the nonlinear response is approximated by
reconstruction points $\{y_i\}_{i=1}^K$. The mean squared error distortion is
\[
D_K = \mathbb{E}\big[(\sigma(X) - Q_K(X))^2\big],
\]
where $\sigma$ denotes the target nonlinearity and $Q_K(X) = y_i$ for $X \in I_i$.

\subsection{Segment-wise Codec Convergence}
\label{app:codec_convergence}

We first establish that the ASNC codec $\mathcal{C}_i(x)$ converges to the optimal
reconstruction point within each segment as training proceeds. This closes the gap
between the learned LIF-based codec and the analytically optimal Lloyd-Max
reconstruction assumed in the distortion bound below.

\begin{lemma}[Segment-wise convergence of ASNC codec]
\label{lem:codec_convergence}
Within segment $i$ covering interval $[a_i, b_i]$, the ASNC codec
$\mathcal{C}_i(x) = \sum_{t=1}^T w_i(t)\,\mathrm{LIF}(\mathrm{quantize}(xs_i +
b_i), \theta_i(t))$ is a universal approximator of piecewise-constant functions on
$[a_i, b_i]$ for sufficiently large $T$. For any target constant $y_i^*$ and any
$\delta > 0$, there exist parameters $\{w_i(t), s_i, b_i, \theta_i(t)\}$ such that
$\sup_{x \in [a_i,b_i]} |\mathcal{C}_i(x) - y_i^*| < \delta$.
\end{lemma}

\begin{proof}
Within segment $i$, the affine transform $x \mapsto xs_i + b_i$ maps $[a_i, b_i]$
to a fixed range $[c_i, d_i]$. After quantization, the LIF unit produces a
deterministic binary spike train for each quantized input level. The output
$\mathcal{C}_i(x) = \sum_t w_i(t)\,\mathrm{spikes}_i(t)$ is a finite linear
combination of indicator functions over quantization levels. For $T$ at least as
large as the number of distinct quantization levels in $[c_i, d_i]$, these indicator
functions span the space of piecewise-constant functions on $[c_i, d_i]$. For a
constant target $y_i^*$, the system $\sum_t w_i(t)\,\mathrm{spikes}_i(t) = y_i^*$
for all spike patterns in the segment is a bounded linear system with a solution.
Standard gradient descent on segment loss $L_i$ converges to this solution.
\end{proof}

\begin{remark}
Lemma~\ref{lem:codec_convergence} implies that at training convergence, each
$\mathcal{C}_i(x)$ approximates the Lloyd-Max optimal reconstruction point $y_i^*$
for segment $i$ to arbitrary precision. The distortion bound below therefore applies
to ASNC at convergence.
\end{remark}

\subsection{Non-uniform Quantization Error Bound}

By Bennett's integral and Lloyd-Max theory~\citep{gersho1992vector}, for any
continuous density $f$ on $\mathbb{R}$ with finite second moment, the high-rate
distortion of the optimal non-uniform quantizer with $K$ intervals is
\[
D_K^{\mathrm{non-uniform}}
= \frac{1}{12}\,K^{-2}
  \left(\int_{-\infty}^{\infty} f(x)^{1/3}\,dx\right)^3
  + o(K^{-2}), \qquad K \to \infty.
\]
For Gaussian $f$, the integral is finite over $\mathbb{R}$, giving asymptotic
constant $C_{\mathrm{non-uniform}} = \tfrac{1}{12}
\bigl(\int_{-\infty}^\infty f(x)^{1/3}\,dx\bigr)^3$.

\subsection{Uniform Partitioning Error Bound}

A uniform scheme partitions $[-R, R]$ into $K$ equal intervals of width $\Delta =
2R/K$. The distortion is
\[
D_K^{\mathrm{uniform}}
= \frac{1}{12}\,\Delta^2 \sum_{i=1}^K \Pr(X \in I_i)
  + 2\int_R^\infty (x-R)^2 f(x)\,dx.
\]
Choosing $R = \Theta(\sqrt{\log K})$ makes the tail term negligible, yielding
\[
D_K^{\mathrm{uniform}}
= \frac{R^2}{3K^2} + o(K^{-2}),
\]
with constant $C_{\mathrm{uniform}} = R^2/3$.

\subsection{Comparison}

By Holder's inequality with exponents $p = 3$ and $q = 3/2$,
\[
\left(\int_{-\infty}^{\infty} f(x)^{1/3}\,dx\right)^3
< R^2 \int_{-\infty}^{\infty} f(x)\,dx = R^2,
\]
where the inequality is strict for any non-constant $f$. Therefore
\[
C_{\mathrm{non-uniform}}
= \frac{1}{12}\left(\int_{-\infty}^{\infty} f(x)^{1/3}\,dx\right)^3
< \frac{R^2}{12}
< C_{\mathrm{uniform}},
\]
and hence $D_K^{\mathrm{ASNC}} < D_K^{\mathrm{SpikeZIP-FT}}$ asymptotically. This
completes the proof of Theorem~\ref{thm:asnc_bound} in the main paper.

\begin{remark}
The extension from the nonnegative half-axis to the full real line is handled
naturally here. For the zero-crossing region where $a_i < 0 < b_i$, the
morphology-aware weighting assigns regional priority $R_{\mathrm{region}} = 3.0$
(Appendix~\ref{app:asnc_details}), ensuring denser partitioning in the region of
highest curvature for SiLU and related activations. This is consistent with the
optimal non-uniform allocation predicted by the Bennett integral, which concentrates
intervals where $f(x)^{1/3}$ is largest.
\end{remark}

\subsection{ASNC Error Localization}
\label{app:error_localization}

\begin{corollary}[ASNC error localization]
\label{cor:asnc_localization}
Let $\epsilon_{\mathrm{ASNC}}$ denote the approximation error at a nonlinear module.
The ASNC output is immediately re-encoded into BSE format
via~\eqref{eq:asnc_output}. By Corollary~\ref{cor:no_diffusion}, all subsequent
linear layers operate on this BSE-encoded output exactly, introducing no additional
error beyond the quantization floor. Therefore $\epsilon_{\mathrm{ASNC}}$ does not
accumulate or amplify: the total network error with $M$ nonlinear modules is bounded
by $M \cdot \epsilon_{\mathrm{ASNC}}$, independent of network depth $L$.
\end{corollary}

\begin{proof}
By construction~\eqref{eq:asnc_output}, the ASNC output $\mathbf{S}^{\mathrm{BSE}}$
is a valid BSE spike train. By Corollary~\ref{cor:no_diffusion}, any linear layer
receiving a valid BSE input produces a valid BSE output with no additional error
beyond the quantization floor. The ASNC approximation error is therefore frozen at
the point of re-encoding and cannot be amplified by subsequent exact linear
operations. Summing over $M$ nonlinear modules gives the stated bound.
\end{proof}

\begin{remark}
This contrasts sharply with prior methods, where per-layer error $\epsilon$ compounds
to $O((1+\epsilon)^L)$ after $L$ layers. Under LASER the bound is
$O(M \cdot \epsilon_{\mathrm{ASNC}})$, linear in $M$ and independent of $L$. For a
typical transformer where $M \ll L$, this represents an exponential-to-linear
improvement in error growth with depth.
\end{remark}


%%%%%%%%%%%%%%%%%%%%%%%%%%%%%%%%%%%%%%%%%%%%%%%%%%%%%%%%%%%%%%%%%%%%%%%%%%%
\section{STE Correctness Under LASER}
\label{app:ste_proof}
%%%%%%%%%%%%%%%%%%%%%%%%%%%%%%%%%%%%%%%%%%%%%%%%%%%%%%%%%%%%%%%%%%%%%%%%%%%

This appendix provides the full proof of Proposition~\ref{prop:ste_main} in the main
paper.

Standard surrogate gradient methods apply STE on top of lossy representations, where
the forward-backward inconsistency is a genuine approximation with no computable error
bound. LASER's forward-path fidelity guarantees allow a precise characterization of
the STE gradient error.

\begin{proposition}[STE correctness under LASER]
\label{prop:ste_correct}
Under the LASER framework, the STE gradient satisfies:
\begin{enumerate}[label=(\roman*),leftmargin=*]
  \item \textbf{Linear layers.} The STE gradient equals the true gradient exactly.
        No approximation is introduced.
  \item \textbf{Nonlinear layers.} The STE gradient error is bounded by
        $\epsilon_{\mathrm{ASNC}} \le D_K^{\mathrm{ASNC}}$
        from Appendix~\ref{app:asnc_proof}.
  \item \textbf{Depth independence.} The total gradient error is bounded by
        $M \cdot \epsilon_{\mathrm{ASNC}}$, independent of network depth $L$.
\end{enumerate}
\end{proposition}

\begin{proof}
\textbf{(i) Linear layers.}
By Corollary~\ref{cor:no_diffusion}, the BSE forward pass through any linear layer
$f(\mathbf{x}) = \mathbf{W}\mathbf{x} + \mathbf{b}$ is functionally identical to
the $N$-bit quantized ANN forward pass. The true gradient with respect to the input
is $\mathbf{W}^\top$. The STE substitutes the identity for the BSE encoder gradient.
Since the encoder is the identity on the integer grid by
Theorem~\ref{thm:bijective}, the chain rule through the linear layer is unaffected
and the STE gradient equals the true gradient exactly.

\textbf{(ii) Nonlinear layers.}
At an ASNC module, the forward pass introduces approximation error
$\epsilon_{\mathrm{ASNC}} \le D_K^{\mathrm{ASNC}}$. The STE passes the upstream
gradient through the module as an identity mapping. The resulting gradient error
relative to the true gradient of $\sigma$ is proportional to
$\epsilon_{\mathrm{ASNC}}$: if the forward output deviates from the true nonlinearity
by $\epsilon_{\mathrm{ASNC}}$, the gradient deviation is bounded by
$\epsilon_{\mathrm{ASNC}} / \delta$, where $\delta$ is the minimum input spacing
within the segment, a quantity controlled by the segment resolution $K$.

\textbf{(iii) Depth independence.}
By Corollary~\ref{cor:asnc_localization}, $\epsilon_{\mathrm{ASNC}}$ does not
propagate or amplify through subsequent linear layers. Each nonlinear module
contributes an independent bounded gradient error; these contributions sum rather than
compound. The total gradient error is therefore $O(M \cdot \epsilon_{\mathrm{ASNC}})$,
independent of network depth $L$.
\end{proof}

\begin{remark}
Proposition~\ref{prop:ste_correct} resolves the long-standing theoretical objection
to STE in SNNs. In prior work, STE is applied on top of lossy representations where
the forward-backward inconsistency is a genuine unbounded approximation with no
computable error guarantee. Under LASER, the inconsistency is zero for linear layers
by Corollary~\ref{cor:no_diffusion} and provably bounded for nonlinear layers by
Theorem~\ref{thm:asnc_bound}. STE is therefore a rigorously justified procedure
within the LASER framework, not a heuristic adopted for lack of a better option.
\end{remark}


%%%%%%%%%%%%%%%%%%%%%%%%%%%%%%%%%%%%%%%%%%%%%%%%%%%%%%%%%%%%%%%%%%%%%%%%%%%
\section{ASNC Formal Specification}
\label{app:asnc_details}
%%%%%%%%%%%%%%%%%%%%%%%%%%%%%%%%%%%%%%%%%%%%%%%%%%%%%%%%%%%%%%%%%%%%%%%%%%%

\subsection{Segmental Codec and BSE Re-encoding}

Each segmental codec $\mathcal{C}_i(x)$ uses a spiking neural structure as defined
in~\eqref{eq:asnc_codec} of the main paper. The decoded analog value $\hat{y}_i$ for
each active segment is immediately re-encoded by the BSE encoder to produce
$\mathbf{S}^{\mathrm{BSE}}_i(t)$ within the same window $T$; the overall output
follows~\eqref{eq:asnc_output}.

\subsection{Adaptive Splitting Algorithm}

\noindent\textbf{Split criterion.}
\begin{equation}
\frac{L_i}{\bar{L}} \cdot \big(1 + \sigma_i^2\big) \cdot \rho_i
\;>\; \tau_{\mathrm{adaptive}},
\end{equation}
where $L_i$ is the per-segment loss, $\bar{L}$ is the global average loss,
$\sigma_i^2$ is the loss variance within the segment, and $\rho_i$ is the activation
proportion.

\noindent\textbf{Adaptive threshold.}
\begin{equation}
\tau_{\mathrm{adaptive}}
\;=\; \tau_{\mathrm{base}} \cdot \alpha_{\mathrm{stag}} \cdot \alpha_{\mathrm{prog}},
\end{equation}
where $\tau_{\mathrm{base}}$ is a base threshold and the $\alpha$ factors correct for
training stagnation and progress:
\begin{equation}
\alpha_{\mathrm{stag}} \;=\;
\begin{cases}
0.2, & \text{if } N_{\mathrm{stag}}/N_{\mathrm{active}} > 0.6,\\
0.4, & \text{if } N_{\mathrm{stag}}/N_{\mathrm{active}} > 0.3,\\
1.0, & \text{otherwise,}
\end{cases}
\qquad
\alpha_{\mathrm{prog}} \;=\;
\begin{cases}
0.4, & \text{if } (L_0-L_t)/L_0 < 0.005,\\
0.6, & \text{if } (L_0-L_t)/L_0 < 0.02,\\
1.0, & \text{otherwise.}
\end{cases}
\end{equation}

\noindent\textbf{Split point selection.}
Candidate split points $x_s$ are selected by maximizing
\begin{equation}
S(x_s)
= \alpha_1\,E_{\mathrm{sep}}(x_s)
+ \alpha_2\,C_{\mathrm{cons}}(x_s)
+ \alpha_3\,B_{\mathrm{bal}}(x_s)
+ \alpha_4\,Y_{\mathrm{sep}}(x_s),
\end{equation}
where $E_{\mathrm{sep}}$ measures error separation, $C_{\mathrm{cons}}$ reflects
internal consistency, $B_{\mathrm{bal}}$ evaluates data balance, and $Y_{\mathrm{sep}}$
captures target separation.

\subsection{Adaptive Weighting Mechanism}

\noindent\textbf{Morphology-aware weights.}
\begin{equation}
w_i = 0.5\,I_{\mathrm{func}}(i) + 0.35\,I_{\mathrm{error}}(i) + 0.15\,D_i,
\end{equation}
where $I_{\mathrm{func}}(i)$ is the functional importance of segment $i$,
$I_{\mathrm{error}}(i)$ quantifies its error profile, and $D_i$ is its data density.

\noindent\textbf{Regional priority $R_{\mathrm{region}}(i)$.}
\begin{equation}
R_{\mathrm{region}}(i) \;=\;
\begin{cases}
3.0, & a_i < 0 \text{ and } b_i > 0 \quad \text{(zero-crossing region)},\\
2.5, & b_i \le 0 \quad \text{(purely negative region)},\\
2.0, & a_i \ge 0 \quad \text{(purely positive region)},\\
1.0, & \text{otherwise.}
\end{cases}
\end{equation}

\end{document}
